\documentclass[11pt,a4paper]{article}

\usepackage[T1]{fontenc}
\usepackage[utf8]{inputenc}
\usepackage[ngerman]{babel}
\usepackage[a4paper,margin=25mm]{geometry}
\usepackage{amsmath,amssymb}
\usepackage{helvet}
\usepackage{sfmath}
\renewcommand{\familydefault}{\sfdefault}
\usepackage{microtype}
\usepackage[hidelinks]{hyperref}
\usepackage{lastpage}
\usepackage{fancyhdr}
\usepackage{parskip}

\pagestyle{fancy}
\fancyhf{}
\fancyfoot[C]{\thepage\ (\pageref{LastPage})}
\renewcommand{\headrulewidth}{0pt}
\renewcommand{\footrulewidth}{0pt}

\hypersetup{
  pdftitle={MoE-Glossar: SDQP / Superdichte QUBLE-Packung},
  pdfauthor={Wolfgang Tscheu}
}

\begin{document}

\vspace*{-1.8cm}

\noindent
\begin{minipage}[t]{0.66\textwidth}
{\LARGE\bfseries SDQP / Superdichte QUBLE-Packung}
\end{minipage}
\hfill
\begin{minipage}[t]{0.30\textwidth}
\raggedleft
{\small MoE-Glossareintrag}\\
{\small Version 261029-1}
\end{minipage}

\vspace{1.4em}

Die Superdichte QUBLE-Packung (SDQP) ist die erste
großräumige Ordnung der QUBLES.\textsuperscript{1}
Sie beschreibt eine frühe Strukturphase extrem hoher
Apeiron-Dichte, in der die einzelnen QUBLES durch
gegenseitige Einspannung zu einer zusammenhängenden
Packung verbunden sind.

Ihre Grundform entsteht schalenweise. Im Zentrum
befindet sich ein kleineres QUBLE, das von zwölf
größeren QUBLES umgeben wird. Die Kontaktlagen
bestimmen die räumliche Anordnung dieses ersten
Verbandes.

Jede weitere Schale passt sich an die bereits
vorhandene Ordnung an. Da die Lücken der
darunterliegenden Schale für die nächste Schicht
nur näherungsweise passen, entsteht eine
geometrische Kompromissstruktur. Die QUBLES
werden dadurch ovoid verformt und stehen unter
permanenter innerer Spannung.

Die SDQP bildet eine dynamisch verspannte
QUBLE-Struktur mit realer Expansion,
Dichtebindung und Schalenbildung.

Die Grenzbedingung der Struktur liegt dort, wo
die innere Spannung, die ovoiden Verformungen
und die Expansionsdifferenzen innerhalb
derselben Packungsordnung ihre Stabilitätsgrenze
erreichen. An dieser Grenze beginnt die
Umordnung zur DQP.\textsuperscript{2}

Im Verlauf der Expansion verändern sich die
Größenverhältnisse und Kontaktlagen der QUBLES.

Das Lückensystem der SDQP bildet ein
durchgängiges Netz gegeneinander versetzter
Zwickel.

Die geometrischen Grundformen von Interstitien
werden im eigenständigen Glossareintrag
\emph{Interstitium} definiert.\textsuperscript{3}

Die spezifische Ausbildung und Verbindung
dieser Zwischenräume ist von der jeweiligen
Packungsordnung abhängig.

Im MoE entspricht die superdichte frühe
Packungsphase den primordialen schwarzen
Löchern: lokal extrem hohe Apeiron-Dichte,
stark verringerte Expansion und schalenweise
gebundene Struktur.

Dieselbe superdichte Packungsordnung liegt
auch späteren schwarzen Löchern zugrunde,
die aus Sternentwicklungen hervorgehen.

Die primordialen schwarzen Löcher unterscheiden
sich davon durch ihre frühe Entstehungsphase
direkt aus der superdichten QUBLE-Ordnung.

Die Emergenzkette lautet:

\[
\text{Apeiron}
\rightarrow
\text{Apeiron-Expansion}
\rightarrow
\text{Quantelung}
\rightarrow
\text{QUBLE}
\rightarrow
\text{SDQP}
\rightarrow
\text{DQP}
\rightarrow
\text{Pfad}
\]

Die SDQP ist damit die erste schalenförmig
aufgebaute, superdichte und dynamisch
verspannte Packungsstufe der Realität.

\vspace{0.9em}
\hrule
\vspace{0.45em}

{\footnotesize
\textsuperscript{1}\ MoE-Glossareintrag
\emph{QUBLE / Quantelung}.\\
\textsuperscript{2}\ MoE-Glossareintrag
\emph{DQP / Dichte QUBLE-Packung}.\\
\textsuperscript{3}\ MoE-Glossareintrag
\emph{Interstitium}.
}

\end{document}
